\documentclass[12pt,a4paper]{article}

\usepackage{fontspec}
\usepackage{xeCJK}
\setCJKmainfont{Hiragino Mincho ProN}
\setCJKsansfont{Hiragino Kaku Gothic ProN}
\setCJKmonofont{Hiragino Kaku Gothic ProN}
\usepackage[margin=2.5cm]{geometry}
\usepackage{amsmath,amssymb}
\usepackage{hyperref}
\usepackage{booktabs}
\usepackage{tcolorbox}
\usepackage{enumitem}
% \usepackage{circuitikz}  % not available

\tcbuselibrary{breakable,skins}

\newtcolorbox{keybox}[1][]{
  colback=black!3, colframe=black!50,
  fonttitle=\bfseries, title={#1},
  breakable, sharp corners
}

\newtcolorbox{warnbox}{
  colback=red!3, colframe=red!40,
  fonttitle=\bfseries, title={注意},
  breakable, sharp corners
}

\newtcolorbox{successbox}{
  colback=green!3, colframe=green!40,
  fonttitle=\bfseries, title={成功したら},
  breakable, sharp corners
}

\setlength{\parskip}{0.8em}
\setlength{\parindent}{0pt}

\begin{document}

\begin{center}
{\LARGE\bfseries 5,100円で宇宙の構造を探る}\\[0.5cm]
{\Large\bfseries トポトロニクスLC回路実験}\\[0.3cm]
{\large --- 高校生のための完全実践ガイド ---}\\[1.5cm]
{\large Wright Brothers}\\[0.3cm]
{2026年4月}
\end{center}

\vspace{0.5cm}

\tableofcontents

\newpage

% ============================================================================
\section{この実験で何がわかるか}
% ============================================================================

この実験は、「宇宙の基本構造は素数でできている」という
仮説の最初の物理的テストだ。

ふつうの電子部品（コンデンサとインダクタ）を並べて回路を作り、
特定のつなぎ方をスイッチで切ると、
回路の振る舞いが劇的に変わる（「エッジ状態」が現れる）。

もしエッジ状態の数が理論の予測どおりなら、
それは「素数が物理法則を支配している」ことの
最初の実験的証拠になる。

\begin{keybox}[実験の概要]
\begin{tabular}{ll}
コスト & 約5,100円 \\
所要時間 & はんだ付け2時間 + 測定1時間 \\
必要な知識 & 中学レベルの電気回路 + FFTの使い方 \\
必要な道具 & はんだごて、PC（イヤホン端子とマイク端子） \\
予測 & スイッチを10個切ると、20個のピークが現れる \\
\end{tabular}
\end{keybox}

% ============================================================================
\section{背景：SSHモデルとは}
% ============================================================================

\subsection{交互につなぐだけで「位相」が変わる}

SSHモデル（Su-Schrieffer-Heeger model）は、
1979年にプラスチックの電気伝導を説明するために作られた理論。

核心は驚くほど単純：

\begin{keybox}[SSHモデルの原理]
素子を一列に並べて、「強い結合」と「弱い結合」を交互に繰り返す。

\textbf{強-弱-強-弱-強-弱...}

たったこれだけで、鎖の両端に「エッジ状態」という特殊な状態が現れる。
これはトポロジカル（位相的）な効果であり、少しの乱れでは消えない。
\end{keybox}

「強い結合」と「弱い結合」の入れ替えは、
数学的には「巻き数」（winding number）という整数で区別される。

\begin{itemize}
\item 弱-強-弱-強... → 巻き数 = 0（ふつう、エッジ状態なし）
\item 強-弱-強-弱... → 巻き数 = 1（トポロジカル、エッジ状態あり）
\end{itemize}

\subsection{「素数ミュート」= 鎖を周期的に切る}

ここからが我々の理論との接続。

鎖の2番目ごとの「強い結合」をスイッチで切断すると、
鎖が複数のブロックに分割される。
各ブロックは独立したSSH鎖になり、
各ブロックの両端にエッジ状態が現れる。

\begin{keybox}[理論の予測（Paper Eから）]
$N$セルの鎖を$p$セルごとに切断すると：
\begin{itemize}
\item ブロック数 = $N/p$
\item エッジ状態の数 = $2 \times N/p$（各ブロックの両端に1つずつ）
\end{itemize}

これは $K_1(\mathbb{Z}) \to K_1(\mathbb{Z}[1/p])$
というK理論の位相遷移に対応する。
\end{keybox}

つまり：スイッチを切る操作 = 「素数$p$のミュート」の古典的アナログ。
エッジ状態の数が予測と一致すれば、K理論的位相遷移が実際に起きた証拠。

% ============================================================================
\section{部品リストと買い方}
% ============================================================================

\subsection{秋葉原での購入リスト}

\begin{table}[h]
\centering
\caption{部品リスト（20セルSSH鎖 + p=2変調）}
\begin{tabular}{@{}llrr@{}}
\toprule
部品 & 規格 & 数量 & 概算価格 \\
\midrule
セラミックコンデンサ & 100 nF (0.1 $\mu$F) & 20個 & 200円 \\
電解コンデンサ & 1 $\mu$F & 19個 & 400円 \\
インダクタ（コイル） & 10 mH & 20個 & 2,000円 \\
トグルスイッチ & 小型ON/OFF & 10個 & 1,000円 \\
ユニバーサル基板 & Cサイズ以上 & 1枚 & 500円 \\
ワイヤー & 単芯（色違い推奨） & 1m×3色 & 300円 \\
はんだ & 鉛フリー & 1巻 & 300円 \\
3.5mmステレオジャック & メス & 2個 & 400円 \\
\midrule
\multicolumn{3}{r}{\textbf{合計}} & \textbf{約5,100円} \\
\bottomrule
\end{tabular}
\end{table}

\textbf{購入場所}：秋月電子通商、千石電商、マルツエレック（秋葉原）、
またはAmazon/MonotaROで通販。

\subsection{追加で必要な道具}

\begin{itemize}
\item はんだごて（温度調整付き推奨、なければ30Wタイプ）
\item PC（Windows/Mac、3.5mmイヤホン出力 + マイク入力）
\item 3.5mmオーディオケーブル 2本
\item FFTソフト: Audacity（無料）
\end{itemize}

\begin{warnbox}
はんだごてでのやけどに注意。換気しながら作業すること。
中学の技術の授業でやったはんだ付けと同じ。不安なら大人に手伝ってもらおう。
\end{warnbox}

% ============================================================================
\section{回路の組み立て}
% ============================================================================

\subsection{回路図の読み方}

SSH型LC回路は以下の構造：

\begin{verbatim}
  回路の構造（1ブロック = 2セル）:

  IN --[Cv]--+--[Cw]--+--[Cv]--+--[SW]--+-- ...
              |         |         |         |
             [L]       [L]       [L]       [L]
              |         |         |         |
  GND --------+---------+---------+---------+-- GND

  Cv = 100 nF（弱い結合、セル内）
  Cw = 1 uF（強い結合、セル間）
  L  = 10 mH（各ノードからグランドへ）
  SW = トグルスイッチ（2セルごとのCwに直列）
\end{verbatim}

\begin{itemize}
\item $C_v = 100\,\mathrm{nF}$（弱い結合 = セル内）
\item $C_w = 1\,\mu\mathrm{F}$（強い結合 = セル間）
\item $L = 10\,\mathrm{mH}$（各ノードからグランドへ）
\item SW = トグルスイッチ（2セルごとの$C_w$に挿入）
\end{itemize}

\subsection{組み立て手順}

\textbf{ステップ1}：ユニバーサル基板に20個のインダクタを一列にはんだ付け。
片側をバス線（グランド）に接続。

\textbf{ステップ2}：インダクタの間に$C_v$（100 nF）を交互にはんだ付け。
これがSSH鎖の「弱い結合」。

\textbf{ステップ3}：残りのインダクタ間に$C_w$（1 $\mu$F）をはんだ付け。
ただし、\textbf{2番目ごとの$C_w$にはスイッチを直列に入れる}。
合計10個のスイッチ。

\textbf{ステップ4}：鎖の両端に3.5mmジャックを接続。
片方が入力（PCのイヤホン出力）、もう片方が出力（PCのマイク入力）。

\textbf{ステップ5}：グランドバスをジャックのグランドに接続。

\begin{keybox}[チェックポイント]
\begin{itemize}
\item[$\square$] インダクタ20個が一列に並んでいる
\item[$\square$] $C_v$（小さいコンデンサ）が奇数番目の間に入っている
\item[$\square$] $C_w$（大きいコンデンサ）が偶数番目の間に入っている
\item[$\square$] $C_w$の2つに1つにスイッチが入っている（10個）
\item[$\square$] 両端にジャックがある
\item[$\square$] 全インダクタがグランドにつながっている
\end{itemize}
\end{keybox}

% ============================================================================
\section{測定方法}
% ============================================================================

\subsection{ソフトウェアの準備}

\textbf{Audacity}（無料音声編集ソフト）をインストール。
\url{https://www.audacityteam.org/}

設定：
\begin{itemize}
\item サンプルレート：44100 Hz
\item 入力：マイク端子
\item 出力：イヤホン端子
\end{itemize}

\subsection{測定1：全スイッチON（標準SSH）}

\textbf{手順}：
\begin{enumerate}
\item 全10個のスイッチをONにする（全結合が生きている状態）
\item PCのイヤホン出力から回路の入力にケーブルを接続
\item 回路の出力からPCのマイク入力にケーブルを接続
\item Audacityの「生成」→「ノイズ」→「ホワイトノイズ」で信号を生成
\item 3秒間録音する
\item 「解析」→「スペクトル表示」でFFTを実行
\item スクリーンショットを保存（これが「基準スペクトル」）
\end{enumerate}

\textbf{期待される結果}：
いくつかの共鳴ピークが見える。鎖の両端にエッジ状態があるので、
低周波に2つの鋭いピークが見えるはず。

\subsection{測定2：スイッチ切断（素数ミュート）}

\textbf{手順}：
\begin{enumerate}
\item 10個のスイッチを全てOFFにする
\item 同じ手順でホワイトノイズを入力、録音、FFT
\item スクリーンショットを保存（これが「ミュート後スペクトル」）
\end{enumerate}

\textbf{理論の予測}：

\begin{keybox}[予測される結果]
スイッチを切ると、鎖が10個のブロックに分割される。
各ブロックの両端にエッジ状態が現れるので：

\textbf{エッジ状態の数 = $2 \times 10 = 20$ 個の追加ピーク}

基準スペクトルにはなかった新しいピークが
約20個出現するはず。
\end{keybox}

\subsection{測定3：部分切断（中間状態）}

\textbf{手順}：スイッチを1個ずつOFFにしていき、
それぞれのスペクトルを記録する。

\begin{itemize}
\item 0個OFF：2ピーク（元のエッジ状態）
\item 1個OFF：4ピーク（+2）
\item 2個OFF：6ピーク（+2）
\item ...
\item 10個OFF：22ピーク（元の2 + 追加20）
\end{itemize}

スイッチを1個切るごとにピークが2つずつ増えれば、
「1回の位相操作でエッジ状態が2つ生まれる」ことの直接証拠。

% ============================================================================
\section{データの解析}
% ============================================================================

\subsection{ピークの数え方}

FFTスペクトルの中で：
\begin{enumerate}
\item ノイズフロアより明らかに高いピークを数える
\item 「全ON」と「全OFF」のスペクトルを重ね合わせて比較する
\item 新しく現れたピークだけをカウントする
\end{enumerate}

\subsection{成功の判定}

\begin{successbox}
\textbf{完全成功}：追加ピーク数 = $2 \times$（OFFにしたスイッチ数）

全スイッチOFFで追加ピークが20個（$\pm$2くらいの誤差は許容）なら、
Paper Eの予測と一致。

\textbf{これは$K_1$トポロジカル相転移の最初の実験的観測であり、
$\Spec(\mathbb{Z})$仮説を支持する最初の物理実験。}
\end{successbox}

\begin{keybox}[部分的成功]
ピーク数は一致しないが、スイッチのON/OFFでスペクトルが
明確に変わる場合 → 「算術的モード選別が物理系で効果を持つ」
ことの証拠にはなる。
\end{keybox}

\begin{warnbox}
ピーク数が全く合わない場合 → 回路の定数（$C_v, C_w, L$の値）が
トポロジカル相の条件を満たしていない可能性がある。
$C_w > C_v$（強い結合 $>$ 弱い結合）が必須条件。
値の比を大きくして（例：$C_w = 10\,\mu$F）やり直してみよう。
\end{warnbox}

% ============================================================================
\section{結果が出たら}
% ============================================================================

\subsection{論文にする}

スペクトルのスクリーンショット（全ON, 全OFF, 部分OFF）と
ピーク数のカウント表があれば、実験ノートとして十分。
GitHubに上げるか、arXivにプレプリントとして投稿できる。

\subsection{次のステップ}

\begin{itemize}
\item \textbf{p=3の検証}：スイッチを3番目ごとに配置して
  $p=3$ミュートを実験。エッジ数 = $2N/3$と一致するか？
\item \textbf{乗法構造}：$p=2$と$p=3$を同時にミュート。
  エッジ数 = $2 \times N/6$（$p=2$と$p=3$の最小公倍数6でブロック化）か？
\item \textbf{SQUID実験への提案}：トポトロニクスの結果を添えて、
  超伝導回路の研究室に共同研究を提案する。
\end{itemize}

% ============================================================================
\section{なぜこの実験が重要か（もう一度）}
% ============================================================================

\begin{keybox}[5,100円の実験の意味]
アインシュタインの相対論は、1919年の日食観測で
「星の光が太陽の重力で曲がる」ことが確認されて
受け入れられた。

この実験は、それの最初の小さなステップ。

「スイッチを切ると予測通りの数のピークが現れる」

もしこれが確認されれば、
「素数と物理法則の間に数学的につながりがある」ことの
最初の実験的証拠になる。

5,100円で、物理学の新しい章が始まるかもしれない。
\end{keybox}

\vfill
\begin{center}
{\small Wright Brothers, 2026}
\end{center}

\end{document}
